\documentclass[11pt]{article}
\usepackage[a4paper,margin=18mm]{geometry}
\usepackage[no-math]{fontspec}
\setmainfont{Latin Modern Roman}
\usepackage{amsmath,amssymb,amsthm,xfrac,xspace,hyperref,graphicx,comment}
\excludecomment{DefTralics}

\providecommand{\bibliofont}{\small}
\newcommand{\ie}{i.e.,\xspace}
\newcommand{\eg}{e.g.,\xspace}
\newcommand{\etc}{etc.\xspace}
\newcommand{\cf}{cf.\xspace}
\newcommand{\resp}{resp.\xspace}
\newcommand{\smaller}{\small}
\newcommand{\Subsection}{\subsection}
\newcommand{\Subsubsection}{\subsubsection}
\newcommand{\skpt}{\smallskip}
\emergencystretch=3em
\let\oldhat\hat\let\oldbar\bar\let\oldtilde\tilde
\numberwithin{equation}{section}\input{author-preamble.tex}
\begin{document}
\begin{center}{\Large Stable item 2037}\end{center}
\paragraph{Source and attribution.}
Andrea D’Agnolo, Masaki Kashiwara, \emph{On a topological counterpart of regularization for holonomic \$\textbackslash{}protect \textbackslash{}mathscr\{D\}\$-modules}, Journal de l’École polytechnique — Mathématiques 8 (2021), publisher version, DOI 10.5802/jep.140. \href{https://jep.centre-mersenne.org/articles/10.5802/jep.140/}{Publisher article}. Authors' text: \href{https://creativecommons.org/licenses/by/4.0/}{CC BY 4.0}.
\paragraph{Editorial problem boundary.}
Prove only Proposition 4.1 below for a morphism of subanalytic bordered spaces, its closed zero locus, a weakly real-constructible enhanced ind-sheaf, and a compact subset of that zero locus. Exact sheafification/strip-object notation, the complete weak-constructible section theorem B.2 with barycentric proof, and the full power-tube cofinality argument are retained. No unrestricted pullback/sheafification commutation or general regularisation adjunction is selected.
\paragraph{Difficulty (editorial): H3.}
The usual enhanced-category and six-functor foundations do not make sheafification commute with closed pullback in this setting. The author proves the needed weak-constructible section formula through explicit barycentric-star neighbourhoods and retractions, then identifies the singular strip-system limit using a power-tube cofinality argument. That new section and cofinality construction, rather than a formal pullback identity, is the remaining nonroutine obligation.
\paragraph{Editorial transcription note.}
Original author language, mathematics and ordering are preserved. Publisher front matter is replaced by this independent article wrapper. Bibliography is recovered from matching cached publisher HTML, including its TeX formulas. No new proof or mathematical repair is supplied. Surrounding original results are retained for conservative context and are not extra selected corpus claims.
\clearpage
\input{author-body.tex}
\begin{thebibliography}{99}
\bibitem{DK16} A. D’Agnolo \& M. Kashiwara - “Riemann-Hilbert correspondence for holonomic D-modules” , Publ. Math. Inst. Hautes Études Sci. 123 (2016), p. 69-197
\bibitem{DK18} A. D’Agnolo \& M. Kashiwara - “A microlocal approach to the enhanced Fourier-Sato transform in dimension one” , Adv. Math. 339 (2018), p. 1-59
\bibitem{DK19} A. D’Agnolo \& M. Kashiwara - “Enhanced perversities” , J. reine angew. Math. 751 (2019), p. 185-241
\bibitem{DK19nu} A. D’Agnolo \& M. Kashiwara - “Enhanced specialization and microlocalization” , 2019
\bibitem{GS14} S. Guillermou \& P. Schapira - “Microlocal theory of sheaves and Tamarkin’s non displaceability theorem” , in Homological mirror symmetry and tropical geometry , Lect. Notes Unione Mat. Ital. , vol. 15 , Springer, Cham, 2014, p. 43-85
\bibitem{IT20} Y. Ito \& K. Takeuchi - “On irregularities of Fourier transforms of regular holonomic $\mathcal{D}$-modules” , Adv. Math. 366 (2020), p. 107093, 62
\bibitem{Kas84} M. Kashiwara - “The Riemann-Hilbert problem for holonomic systems” , Publ. RIMS, Kyoto Univ. 20 (1984) no. 2, p. 319-365
\bibitem{Kas16} M. Kashiwara - “Riemann-Hilbert correspondence for irregular holonomic $\mathcal{D}$-modules” , Japan. J. Math. 11 (2016) no. 1, p. 113-149
\bibitem{KS90} M. Kashiwara \& P. Schapira - Sheaves on manifolds , Grundlehren Math. Wiss. , vol. 292 , Springer-Verlag, Berlin, Heidelberg, 1990
\bibitem{KS01} M. Kashiwara \& P. Schapira - Ind-sheaves , Astérisque , vol. 271 , Société Mathématique de France, Paris, 2001
\bibitem{KS03} M. Kashiwara \& P. Schapira - “Microlocal study of ind-sheaves. I. Micro-support and regularity” , in Autour de l’analyse microlocale , Astérisque , vol. 284 , Société Mathématique de France, Paris, 2003, p. 143-164
\bibitem{KS16} M. Kashiwara \& P. Schapira - Regular and irregular holonomic D-modules , London Math. Soc. Lect. Note Series , vol. 433 , Cambridge University Press, Cambridge, 2016
\bibitem{Tam08} D. Tamarkin - “Microlocal condition for non-displaceability” , in Algebraic and analytic microlocal analysis , Springer Proc. Math. Stat. , vol. 269 , Springer, Cham, 2018, p. 99-223
\end{thebibliography}
\end{document}
